\documentclass[14pt]{extarticle}

% ---------- PACKAGES ----------
\usepackage{amsmath, amssymb}
\usepackage{geometry}
\usepackage{setspace}
\usepackage{titlesec}
\usepackage{xcolor}
\usepackage{helvet}
\usepackage{enumitem}   % for clean enumerate labels
\usepackage{ifthen}     % for simple toggles

\makeatletter
\IfFileExists{../common-things/reflectionpage.tex}{%
  \input{../common-things/reflectionpage.tex}%
}{%
  \IfFileExists{common-things/reflectionpage.tex}{%
    \input{common-things/reflectionpage.tex}%
  }{%
    \PackageError{\jobname}{Missing reflectionpage.tex file}{%
      Place reflectionpage.tex inside a common-things/ directory next to this unit\@.%
    }%
  }%
}
\makeatother

% ---------- PAGE SETUP ----------
\geometry{margin=1in}
\setstretch{1.5}
\setlength{\parindent}{0pt}
\setlength{\parskip}{0.75em}
\setlist{itemsep=0.75\baselineskip, topsep=0.5\baselineskip}
\titleformat{\section}{\normalfont\Large\bfseries}{\thesection}{1em}{}
\titleformat{\subsection}{\normalfont\large\bfseries}{\thesubsection}{1em}{}
\renewcommand{\familydefault}{\sfdefault}
\renewcommand{\emph}[1]{\textbf{#1}}

% ---------- TOGGLES ----------
\newboolean{showsolutions}
\setboolean{showsolutions}{true} % change to false for student version

% ---------- LIGHTWEIGHT MACROS ----------
\newcommand{\UnitTitle}{Unit X: <<Unit Title>>}
\newcommand{\TopicTitle}{Topic Y: <<Topic Title>>}

% Metadata (purely informational, helps future automation)
% Example: \TopicMeta{Heart of Algebra}{No-Calculator}{Core → Mixed → Challenge}
\newcommand{\TopicMeta}[3]{%
  \textcolor{gray}{\footnotesize \textbf{SAT Domain:} #1 \quad \textbf{Calculator:} #2 \quad \textbf{Progression:} #3}
}

% Inline problem tags (optional; shows small gray tag)
% Example: \qtag{HoA, no-calc, medium}
\newcommand{\qtag}[1]{\textcolor{gray}{\footnotesize [#1]}}

% Example block (keeps your voice + formatting)
\newenvironment{Example}[1]{%
  \subsection*{Example #1}%
}{\vspace{0.25em}}

% Practice part block (Parts A, B, C…)
\newenvironment{PracticePart}[1]{%
  \subsection*{Part #1}%
  \begin{enumerate}
}{%
  \end{enumerate}
}

% “Speed Check” (optional quick 3 problems, no-calc warmup)
\newenvironment{SpeedCheck}{%
  \subsection*{Speed Check (3 quick items)}
  \begin{enumerate}
}{%
  \end{enumerate}
}

% Solution item macro (keeps solution formatting uniform)
% Usage inside solutions: \sol{1}{x = 5}
\newcommand{\sol}[2]{\item[\textbf{(#1)}] #2}

% ---------- DOCUMENT ----------
\begin{document}
\raggedright
\pagenumbering{gobble}

\begin{center}
  \LARGE \textbf{\UnitTitle}\\[6pt]
  \Large \textbf{\TopicTitle}\\[4pt]
  \TopicMeta{<<Domain>>}{<<Calc/No-Calc>>}{Core \textrightarrow\ Mixed \textrightarrow\ Challenge}
\end{center}

\vspace{0.75em}

% =========================
% MINI-LESSON
% =========================
\section*{Concept Summary}
<<2–5 tight paragraphs that define the central idea. Keep formal notation and 1 “why it works” note.>>
\[
\text{(Include one canonical form here, e.g., } ax+b=c \text{ or system } \begin{cases} \cdot \\ \cdot \end{cases})
\]

\section*{Core Skills}
\begin{itemize}
  \item <<Skill 1: action verb + precise behavior>>
  \item <<Skill 2>>
  \item <<Skill 3>>
\end{itemize}

\begin{Example}{1}
Solve for \(x\): \(\;<<expression>>\) \qtag{<<tags>>}

\textbf{Step 1:} <<What and why>> \\
\textbf{Step 2:} <<What and why>>
\[
<<clean working>>
\]
\textbf{Check:} <<Substitute/interpretation>> \quad \(\checkmark\)

\textbf{Final Answer:} \(\boxed{<<value/statement>>}\)
\end{Example}

\begin{Example}{2}
<<Second example with a different twist (fractions, parentheses, or interpretation).>>
\end{Example}

\section*{Key Takeaways}
\begin{itemize}
  \item <<Principle 1 (actionable)>>
  \item <<Principle 2 (pitfall to avoid)>>
  \item <<Principle 3 (verification habit)>>
\end{itemize}

\newpage

% =========================
% PRACTICE SET
% =========================
\section*{Practice Questions: <<Topic Title>>}

% Optional: quick warmup
% \begin{SpeedCheck}
%   \item <<No-calc, 15–30 sec>>
%   \item <<No-calc, 15–30 sec>>
%   \item <<No-calc, 15–30 sec>>
% \end{SpeedCheck}

\begin{PracticePart}{A: Core (one-step / direct use)}
  \item <<problem>> \qtag{<<tags>>}
  \item <<problem>> \qtag{<<tags>>}
  \item <<problem>>
  \item <<problem>>
  \item <<problem>>
\end{PracticePart}

\begin{PracticePart}{B: Two-Step / Structured Reasoning}
  \setcounter{enumi}{5}
  \item <<problem>>
  \item <<problem>>
  \item <<problem>>
  \item <<problem>>
  \item <<problem>>
\end{PracticePart}

\begin{PracticePart}{C: Variables on Both Sides / Structure Recognition}
  \setcounter{enumi}{10}
  \item <<problem>>
  \item <<problem>>
  \item <<problem>>
  \item <<problem>>
  \item <<problem>>
\end{PracticePart}

\begin{PracticePart}{D: Parentheses / Fractions / Edge Cases}
  \setcounter{enumi}{15}
  \item <<problem>>
  \item <<problem>>
  \item <<problem>>
  \item <<problem>>
  \item <<problem>>
\end{PracticePart}

\newpage

% =========================
% SOLUTIONS
% =========================
\ifthenelse{\boolean{showsolutions}}{%
  \section*{Answer Key and Solutions}
  \begin{enumerate}
    \sol{1}{<<solution>>}
    \sol{2}{<<solution>>}
    \sol{3}{<<solution>>}
    \sol{4}{<<solution>>}
    \sol{5}{<<solution>>}
    % Continue numbering or use \setcounter as needed
  \end{enumerate}
}{%
  % Leave blank when solutions are hidden
}

\newpage

\ReflectionPage{<<Next Topic Title>>}{https://www.scotchegg.co}

\end{document}
